\chapter{Lagrangian Mechanics}
\label{chap:lagrangian}

\section{Mechanical Data}
\label{sec:lagrangian-mechanical-data}

A mechanics problem begins by naming its space of possible configurations. Once
that space is named, the rest of the theory can be formulated independently of
any particular coordinate chart.

A pendulum moves on a circle, not on a line. A double pendulum lives on a
torus. A rigid body lives on \(\SO(3)\), not in a box of nine independent
matrix entries. Choosing the configuration space correctly is therefore not a
matter of taste; it is the first physical idealization. If the chosen space
contains redundant coordinates, constraints must later repair that redundancy.
If the chosen space already reflects the physical degrees of freedom, the
variational principle can act directly on the geometry of the problem.

\begin{definition}[Finite-dimensional mechanical system]
A regular finite-dimensional Lagrangian mechanical system consists of:
\begin{enumerate}
\item a smooth \(n\)-dimensional configuration manifold \(Q\);
\item a time interval \(I=[t_0,t_1]\);
\item a Lagrangian \(L:TQ\times I\to \R\);
\item a class of admissible curves \(q:I\to Q\), possibly satisfying constraints.
\end{enumerate}
The action of an admissible curve is
\[
  S[q]=\int_{t_0}^{t_1} L(q(t),\dot q(t),t)\,\dd t.
\]
\end{definition}

\begin{setting}[Coordinate patch for derivations]
Derivations in this chapter are written in one coordinate chart
\((q^1,\ldots,q^n)\). The equations obtained are tensorial or intrinsic when
the objects used to build them are intrinsic. Coordinates are a calculation
device, not an additional physical structure.
\end{setting}

The notation \(q(t)\) should be read as a point of \(Q\) depending on time, not
as a column vector until a chart has been chosen.  In a chart the same motion is
represented by component functions \(q^i(t)\), and the velocity is represented
by \(\dot q^i(t)\).  The action is therefore not a functional of arbitrary
Euclidean functions in the first instance; it is a functional on curves in the
configuration space.  This distinction becomes essential for rigid bodies,
constraints, and later for continua, where a careless coordinate choice can
hide the actual degrees of freedom.

\paragraph{Roadmap.}
The chapter has three jobs. First, it derives the Euler-Lagrange equations from
Hamilton's principle, including boundary terms, constraints, and Noether
charges. Second, it treats Kepler motion and the rigid body as examples in
which those ideas control a complete calculation. Third, it prepares the
boundary momentum and energy concepts that become phase-space coordinates in
Hamiltonian mechanics. The
variational principle separates coordinate choices from geometric structure and
places forces, constraints, and conserved quantities in a common framework.

The primary object here is a path in configuration space, and the variational
principle selects the physical path. The phase-space viewpoint is developed in
Chapter~\ref{chap:hamiltonian-phase-space}. In that order, momenta and
symplectic forms appear as consequences of the endpoint variation and the
Legendre transform.

\paragraph{Derivation checkpoints.}
Whenever this chapter differentiates an action, check three things before
following the algebra: which variables are varied, which endpoint conditions
are imposed, and which boundary terms survive. Whenever it invokes a symmetry,
check whether the symmetry acts on configuration space and whether the action
is invariant up to an allowed boundary term. These checks are small, but they
prevent the most common mistakes: using an inadmissible virtual displacement or
calling a quantity conserved before verifying the relevant symmetry condition.

\section{Variations And The Euler-Lagrange Equations}
\label{sec:euler-lagrange}

What does it mean for the action to be stationary? It means that the first
derivative of \(S[q_\eps]\) vanishes for every admissible infinitesimal change
of the path.

\paragraph{Local notation.}
Throughout this section \(q(t)\) is a path in the configuration manifold \(Q\),
and \(q^i(t)\) are its components in a chosen local coordinate chart.  The
index \(i=1,\ldots,n\) labels configuration coordinates, and repeated upper and
lower \(i\)-indices are summed.  The symbol \(\eps\) labels a family of nearby
paths; it is not a physical time or a small parameter in the Lagrangian itself.

\begin{postulate}[Hamilton's principle for prescribed endpoints]
For the basic Lagrangian formalism used in this section, the physical path with
given endpoint configurations \(q(t_0)\) and \(q(t_1)\) is a stationary point
of the action with respect to all smooth variations that keep those endpoints
fixed.  This is a postulate of the variational description.  It is not a
claim that the physical system has an intrinsic tendency to keep endpoints
fixed; rather, it defines the boundary-value problem whose stationary curves
are being studied.
\end{postulate}

\begin{remark}[Stationary, not necessarily least]
Hamilton's principle, formulated in 1834--1835, is a critical-point principle
for the time integral of \(L=T-V\). The older phrase ``principle of least
action'' refers to a different functional: the action at fixed energy, due to
Maupertuis (1744) and Euler. In modern Lagrangian mechanics the action need not
be a minimum; it may be a saddle or a maximum for some classes of variations.
The catenoid soap film is the familiar variational counterexample: a stationary
surface need not be the minimizing surface once the boundary data are large
enough. The word stationary is therefore the accurate word: the first variation
vanishes.
\end{remark}

Let \(q_\eps(t)\) be a smooth one-parameter family of curves with
\(q_0(t)=q(t)\). The variation vector field along \(q\) is
\[
  \eta(t)=\left.\frac{\partial q_\eps(t)}{\partial \eps}\right|_{\eps=0},
  \qquad
  \eta^i(t)=\left.\frac{\partial q_\eps^i(t)}{\partial \eps}\right|_{\eps=0}.
\]
For the fixed-endpoint variational postulate we impose
\[
  \eta(t_0)=0,\qquad \eta(t_1)=0.
\]
The variational calculation below has one essential idea. The path \(q(t)\) is
a critical point of the action when every infinitesimal deformation of the path
changes the action only at second order. Differentiating the action produces a
term proportional to \(\dot\eta\); integration by parts moves the derivative
onto the momentum, and the endpoint condition kills the boundary term. What is
left must vanish for every interior virtual displacement \(\eta\), which forces
the Euler-Lagrange equations pointwise in time.

\begin{figure}[htbp]
\centering
\begin{tikzpicture}[scale=1.0, >=Latex]
  \draw[->] (-0.3,0) -- (7.45,0) node[right] {\(t\)};
  \draw[->] (0,-0.2) -- (0,3.0) node[above] {configuration coordinate};
  \draw[densely dotted] (0.5,0) -- (0.5,2.75);
  \draw[densely dotted] (6.5,0) -- (6.5,2.75);
  \draw[thick, blue] plot[smooth] coordinates {(0.5,0.7) (1.8,1.25) (3.2,1.0) (4.7,2.0) (6.5,2.35)};
  \draw[thick, red, dashed] plot[smooth] coordinates {(0.5,0.7) (1.8,1.55) (3.2,1.25) (4.7,1.75) (6.5,2.35)};
  \draw[->, red] (3.2,1.0) -- (3.2,1.25);
  \node[red, fill=white, inner sep=1pt] at (3.58,1.18) {\(\eta(t)\)};
  \draw[->, red!70!black] (4.7,2.0) -- (4.7,1.75);
  \fill (0.5,0.7) circle (2pt) node[below right, fill=white, inner sep=1pt] {\(q(t_0)\)};
  \fill (6.5,2.35) circle (2pt);
  \node[below left, fill=white, inner sep=1pt] at (6.44,2.28) {\(q(t_1)\)};
  \node[below] at (0.5,0) {\(t_0\)};
  \node[below] at (6.5,0) {\(t_1\)};
  \node[blue, fill=white, inner sep=1pt] at (2.3,0.55) {\(q(t)\)};
  \node[red, align=left, fill=white, inner sep=1.5pt] (varlabel) at (4.55,2.68)
    {\(q_\eps(t)=q(t)+\eps\eta(t)+O(\eps^2)\)};
  \draw[red!70!black, ->] (varlabel.south west) -- (4.55,1.83);
\end{tikzpicture}
\caption{A fixed-endpoint variation changes the path in the interior while
holding the initial and final configurations fixed. The infinitesimal change is
the vector field \(\eta(t)=\delta q(t)\) along the path.}
\label{fig:fixed-endpoint-variation}
\end{figure}

\begin{derivation}[Euler-Lagrange equations]
The first variation of the action is
\[
\begin{aligned}
  \delta S
  &=
  \left.\frac{\dd}{\dd\eps}\right|_{\eps=0}
  \int_{t_0}^{t_1}L(q_\eps,\dot q_\eps,t)\,\dd t\\
  &=
  \int_{t_0}^{t_1}
  \left(
    \frac{\partial L}{\partial q^i}\eta^i
    +
    \frac{\partial L}{\partial \dot q^i}\dot\eta^i
  \right)\dd t.
\end{aligned}
\]
Integrate the second term by parts:
\[
  \int_{t_0}^{t_1}
  \frac{\partial L}{\partial \dot q^i}\dot\eta^i\,\dd t
  =
  \left[
    \frac{\partial L}{\partial \dot q^i}\eta^i
  \right]_{t_0}^{t_1}
  -
  \int_{t_0}^{t_1}
  \frac{\dd}{\dd t}\left(\frac{\partial L}{\partial \dot q^i}\right)\eta^i
  \,\dd t.
\]
The boundary term vanishes because \(\eta(t_0)=\eta(t_1)=0\). Therefore
\[
  \delta S
  =
  \int_{t_0}^{t_1}
  \left[
    \frac{\partial L}{\partial q^i}
    -
    \frac{\dd}{\dd t}\left(\frac{\partial L}{\partial \dot q^i}\right)
  \right]\eta^i\,\dd t.
\]
The fundamental lemma of the calculus of variations says that if this integral
vanishes for all compactly supported variations \(\eta\), then the coefficient
of every \(\eta^i\) vanishes. Hence
\[
  \boxed{
  \frac{\dd}{\dd t}\left(\frac{\partial L}{\partial \dot q^i}\right)
  -
  \frac{\partial L}{\partial q^i}=0
  }.
\]
\end{derivation}

\begin{proof}[Fundamental lemma used above]
Under the standing smoothness hypothesis on \(L\), \(q\), and the variation,
the bracketed coefficient in the variation formula is continuous. The lemma
therefore applies with \(f\) equal to that coefficient. Let \(f(t)\) be
continuous on \((t_0,t_1)\), and suppose
\[
  \int_{t_0}^{t_1} f(t)\eta(t)\,\dd t=0
\]
for every smooth compactly supported test function \(\eta\). If \(f(t_*)>0\)
at some interior point, continuity gives an interval \([a,b]\) on which
\(f\ge f(t_*)/2\). Choose a nonnegative smooth bump \(\eta\) supported in
\([a,b]\) and not identically zero. Then the integral is positive, a
contradiction. Applying the same argument to \(-\eta\) rules out \(f(t_*)<0\).
Thus \(f=0\) pointwise. Applying this to each component of the bracketed
coefficient gives the Euler-Lagrange equations.
\end{proof}

The Euler-Lagrange equation is local in time because the variations are
arbitrary in the interior of the interval. A single global statement,
\(\delta S=0\), becomes a pointwise differential equation because
\(\eta(t)\) can be chosen independently near each time. Boundary conditions and
endpoint terms are part of the variational data: they decide which variations
are admissible and which physical quantities appear at the endpoints.

For a regular Lagrangian, the velocity Hessian
\[
  W_{ij}(q,\dot q,t)
  =
  \frac{\partial^2L}{\partial\dot q^i\partial\dot q^j}
\]
is nonsingular. The Euler-Lagrange equations can then be solved locally for
\(\ddot q^i\) as smooth functions of \((q,\dot q,t)\). Standard ODE existence
and uniqueness therefore gives a unique local solution for each initial state
\((q(t_0),\dot q(t_0))\), as long as the solution remains in the chosen chart
and no singularity is reached.

\section{Boundary Terms And Forces}
\label{sec:boundary-forces}

Integration by parts also defines momentum. Without imposing fixed endpoints,
the boundary contribution is
\[
  \left[p_i\eta^i\right]_{t_0}^{t_1},
  \qquad
  p_i=\frac{\partial L}{\partial \dot q^i}.
\]
Thus \(p_i\) is the coefficient of endpoint displacement in the first variation
of the action. This boundary-term definition reduces to ``mass times velocity''
only in special coordinates for special Lagrangians.

\paragraph{Local notation.}
Here \(p_i\) is the covector conjugate to \(q^i\), while \(\eta^i=\delta q^i\)
is an allowed virtual displacement.  The product \(p_i\eta^i\) is a scalar
endpoint contribution to the action variation.  The generalized force \(Q_i\)
introduced below is also a covector, so \(Q_i\eta^i\) is virtual work.

If nonconservative generalized forces \(Q_i(q,\dot q,t)\) are present, the
Lagrange-d'Alembert principle adds virtual work to Hamilton's principle. This
is the variational form of d'Alembert's 1743 virtual-work principle, later
absorbed into Lagrange's analytic mechanics. The admissible virtual
displacements are the same ones used in the variational problem, and the
external force does virtual work \(Q_i\eta^i\). Thus
\[
  \delta S+\int_{t_0}^{t_1}Q_i\eta^i\,\dd t=0,
\]
which gives
\[
  \frac{\dd}{\dd t}\left(\frac{\partial L}{\partial \dot q^i}\right)
  -
  \frac{\partial L}{\partial q^i}
  =
  Q_i.
\]
This formula is also the bridge to continua, where the virtual work is supplied
by tractions, pressure forces, or dissipative stresses.

The generalized force \(Q_i\) is a covector along the path, not necessarily a
Euclidean force vector. In Cartesian coordinates for a particle the distinction
is invisible, but in generalized coordinates \(Q_i\eta^i\) is the virtual work
done by the force under the virtual displacement \(\eta\). This is why the
same variational formula can describe torques, constraint reactions, pressure
work, and elastic boundary tractions once the correct configuration space has
been named.

\section{Natural Lagrangians And The Metric Form}
\label{sec:natural-lagrangians}

Many mechanical systems have a Lagrangian of the form
\[
  L(q,\dot q)=T(q,\dot q)-V(q)
  =
  \frac{1}{2}g_{ij}(q)\dot q^i\dot q^j - V(q),
\]
where \(g_{ij}\) is a positive-definite mass metric on \(Q\), and \(V\) is
potential energy.

\paragraph{Local notation.}
The matrix \(g_{ij}(q)\) lowers indices and \(g^{ij}(q)\) denotes its inverse.
The symbols \(\Gamma^\ell_{jk}\) below are Christoffel symbols of this mass
metric.  They are connection coefficients in the chosen coordinates, so they
do not transform as tensor components by themselves.

\begin{remark}[Terminology]
Here ``natural'' is conventional terminology: the kinetic energy is quadratic
in velocity and the remaining term is a potential energy on configuration
space.
\end{remark}

The word ``metric'' here has physical content. It tells us how much kinetic
energy a given velocity has, and it may depend on configuration. For a point
particle in Cartesian coordinates \(g_{ij}=m\delta_{ij}\). For a
pendulum, double pendulum, or rigid body written in angles, the same kinetic
energy becomes a metric on the configuration space. Christoffel symbols can
appear for two distinct reasons that should not be confused. In flat Euclidean
space written in curvilinear coordinates they are coordinate inertial terms; on
a genuinely curved configuration space they also encode the curvature of the
mass metric. In both cases they are not new applied forces.

\begin{derivation}[Geodesic-plus-force equation]
For the natural Lagrangian,
\[
  \frac{\partial L}{\partial \dot q^i}=g_{ij}\dot q^j,
\]
so
\[
  \frac{\dd}{\dd t}\frac{\partial L}{\partial \dot q^i}
  =
  g_{ij}\ddot q^j
  +
  \frac{\partial g_{ij}}{\partial q^k}\dot q^k\dot q^j.
\]
Also
\[
  \frac{\partial L}{\partial q^i}
  =
  \frac{1}{2}\frac{\partial g_{jk}}{\partial q^i}\dot q^j\dot q^k
  -
  \frac{\partial V}{\partial q^i}.
\]
The Euler-Lagrange equation becomes
\[
  g_{ij}\ddot q^j
  +
  \left(
    \frac{\partial g_{ij}}{\partial q^k}
    -
    \frac{1}{2}\frac{\partial g_{jk}}{\partial q^i}
  \right)\dot q^k\dot q^j
  =
  -\frac{\partial V}{\partial q^i}.
\]
After symmetrizing the velocity terms and raising an index with \(g^{\ell i}\),
\[
  \boxed{
  \ddot q^\ell+\Gamma^\ell_{jk}\dot q^j\dot q^k
  =
  -g^{\ell i}\frac{\partial V}{\partial q^i}
  },
\]
where
\[
  \Gamma^\ell_{jk}
  =
  \frac{1}{2}g^{\ell i}
  \left(
    \frac{\partial g_{ij}}{\partial q^k}
    +
    \frac{\partial g_{ik}}{\partial q^j}
    -
    \frac{\partial g_{jk}}{\partial q^i}
  \right)
\]
are the Christoffel symbols of the mass metric.
\end{derivation}

Thus even ``force-free'' motion on a curved configuration space need not look
like \(\ddot q=0\) in coordinates. The Christoffel terms are inertial terms
created by the geometry of \(Q\) and the chosen coordinates. The invariant
content is that the covariant acceleration equals the metric gradient of the
potential:
\[
  \nabla_{\dot q}\dot q=-\operatorname{grad}_g V.
\]
This is the coordinate-free form of the boxed equation.

\section{Holonomic Constraints}
\label{sec:holonomic-constraints}

A holonomic constraint is a condition
\[
  \Phi^a(q,t)=0,\qquad a=1,\ldots,k,
\]
that cuts out an admissible submanifold of configuration space. If one works in
the ambient coordinates \(q^i\), admissible variations must preserve the
constraint to first order:
\[
  \frac{\partial\Phi^a}{\partial q^i}\eta^i=0.
\]

\paragraph{Local notation.}
The label \(a=1,\ldots,k\) counts independent constraint equations, while
\(i=1,\ldots,n\) counts ambient configuration coordinates.  The multipliers
\(\lambda_a(t)\) introduced below are unknown functions of time whose role is
to enforce \(\Phi^a(q(t),t)=0\) along the motion.

The Euler-Lagrange residual therefore has to vanish only on vectors tangent to
the constraint surface. The multiplier rule, in the form Lagrange used in
\emph{Mecanique analytique} (1788), says that a covector
vanishing on this tangent space is a linear combination of the constraint
gradients. The linear-algebra statement is this: when the \(k\) one-forms
\(\dd\Phi^a\) are independent, the tangent space is
\[
  T_qC=\{\eta:\dd\Phi^a(\eta)=0,\ a=1,\ldots,k\},
\]
and its annihilator is
\[
  (T_qC)^0=\operatorname{span}\{\dd\Phi^1,\ldots,\dd\Phi^k\}.
\]
Indeed, the map \(A:T_qQ\to\R^k\), \(A(\eta)=(\dd\Phi^1(\eta),\ldots,
\dd\Phi^k(\eta))\), has rank \(k\), so its dual map
\(A^*(\lambda)=\lambda_a\dd\Phi^a\) has image equal to the annihilator of
\(\ker A=T_qC\). This is the finite-dimensional fact behind the multiplier
rule.
Thus constraint regularity means the gradients of the constraints have
constant rank \(k\) on the constraint surface. Under this regularity hypothesis,
the constrained equations take the multiplier form
\[
  \frac{\dd}{\dd t}\frac{\partial L}{\partial \dot q^i}
  -
  \frac{\partial L}{\partial q^i}
  =
  \lambda_a\frac{\partial \Phi^a}{\partial q^i}.
\]
The multipliers \(\lambda_a\) are not arbitrary forces. They are whatever
constraint reactions are needed to keep the motion on \(\Phi^a=0\). If instead
one parametrizes the constraint surface directly, the \(\lambda_a\) disappear
and the same physics is expressed with fewer coordinates.

\section{Symmetry And Noether's Theorem}
\label{sec:noether}

Continuous symmetries of the action produce conserved quantities. In the form
used in this section, Noether's 1918 theorem applies to transformations of
configuration variables that leave the time coordinate itself fixed. The
infinitesimal displacement of \(q\) may nevertheless depend explicitly on time:
a Galilean boost changes a path by \(\delta x=vt\), even though it does not
change the time label \(t\). A symmetry leaves the action unchanged, or changes
it only by an endpoint term. The conserved quantity is the momentum paired with
the infinitesimal displacement, corrected by the endpoint term when one is
present. More general versions of Noether's theorem also allow transformations
of time and fields; those require the extended formalism introduced later when
the Hamiltonian and the action are treated on phase space.

\paragraph{Local notation.}
In this section \(G\) is a Lie group acting on configuration space \(Q\),
\(\mathfrak g\) is its Lie algebra, and \(\xi\in\mathfrak g\) is an
infinitesimal symmetry parameter.  The vector field \(\xi_Q\) is the induced
motion of points of \(Q\).  When a symmetry is explicitly time-dependent, we
write this displacement as \(\xi_Q(q,t)\); it is still vertical, meaning that
it changes \(q\) at the same time \(t\).  The scalar \(B_\xi\) records a
possible endpoint term, and \(J_\xi\) is the corresponding conserved Noether
charge.

\begin{definition}[Infinitesimal generator]
Let a Lie group \(G\) act on \(Q\). For \(\xi\in\mathfrak g\), the infinitesimal
generator is the vector field
\[
  \xi_Q(q)=\left.\frac{\dd}{\dd s}\right|_{s=0}\exp(s\xi)\cdot q.
\]
In coordinates, write \(\xi_Q=\xi_Q^i(q)\partial/\partial q^i\).
The later proof also applies to explicitly time-dependent vertical
displacements \(\xi_Q(q,t)\), which should be read as infinitesimal
transformations of paths rather than as time-independent actions on \(Q\).
\end{definition}

\begin{definition}[Lagrangian symmetry]
The action of \(G\) is a symmetry if
\[
  L(g\cdot q,\,T_qg\cdot \dot q,\,t)=L(q,\dot q,t)
\]
for all \(g\in G\), \(q\in Q\), and \(\dot q\in T_qQ\).
\end{definition}

Some important symmetries preserve the action only up to an endpoint term. If
the infinitesimal variation of the Lagrangian has the form
\[
  \delta_\xi L=\frac{\dd}{\dd t}B_\xi(q,t),
\]
then the action changes by \(B_\xi(t_1)-B_\xi(t_0)\), so fixed-endpoint
variations still give the same equations of motion. The conserved Noether
quantity is then
\[
  J_\xi=p_i\xi_Q^i-B_\xi.
\]
Galilean boosts of a free particle are the standard example: the Lagrangian is
not strictly invariant, but its change is a total time derivative, and the
boost charge includes the corresponding boundary term.  In that example
\(\xi_Q=t\,\partial_x\) depends on time, which is why the vertical
time-dependent formulation is needed.

\begin{proposition}[Noether theorem]
For every \(\xi\in\mathfrak g\), the quantity
\[
  J_\xi(q,\dot q)
  =
  \frac{\partial L}{\partial \dot q^i}\xi_Q^i(q,t)
  -B_\xi(q,t)
  =
  p_i\xi_Q^i(q,t)-B_\xi(q,t)
\]
is conserved along solutions of the Euler-Lagrange equations whenever
\(\delta_\xi L=\dd B_\xi/\dd t\). For a strict symmetry, \(B_\xi=0\).
\end{proposition}

\begin{proof}
Differentiate the symmetry condition along \(g=\exp(s\xi)\) at \(s=0\):
\[
  \frac{\dd B_\xi}{\dd t}=
  \frac{\partial L}{\partial q^i}\xi_Q^i
  +
  \frac{\partial L}{\partial \dot q^i}
  \frac{\dd}{\dd t}\xi_Q^i(q(t),t).
\]
Along an Euler-Lagrange solution,
\[
  \frac{\partial L}{\partial q^i}
  =
  \frac{\dd}{\dd t}\frac{\partial L}{\partial \dot q^i}.
\]
Substitute this into the previous equation:
\[
  \frac{\dd B_\xi}{\dd t}=
  \frac{\dd}{\dd t}\left(\frac{\partial L}{\partial \dot q^i}\right)\xi_Q^i
  +
  \frac{\partial L}{\partial \dot q^i}
  \frac{\dd}{\dd t}\xi_Q^i
  =
  \frac{\dd}{\dd t}
  \left(
    \frac{\partial L}{\partial \dot q^i}\xi_Q^i
  \right).
\]
Therefore \(\dd(p_i\xi_Q^i-B_\xi)/\dd t=0\).
\end{proof}

\begin{example}[Galilean boost as a boundary-term symmetry]
For a free particle on the line,
\[
  L=\frac12m\dot x^2.
\]
An infinitesimal boost with parameter \(v\) changes the path by
\(\delta x=vt\) and \(\delta\dot x=v\). Thus
\[
  \delta L=m\dot x\,v=\frac{\dd}{\dd t}(mvx).
\]
Here \(B=mvx\), and the Noether charge is
\[
  J=p\,vt-mvx=mv(\dot x\,t-x).
\]
Removing the arbitrary boost parameter gives \(m(\dot x\,t-x)\); its negative
\(m(x-\dot x\,t)\) is the conventional conserved center-of-mass quantity for a
free particle. The Lagrangian was not invariant as a function, but the action
changed only by an endpoint term.
\end{example}

\begin{example}[Translation and rotation]
For a particle in \(\R^3\) with \(L=\frac12m\norm{\dot x}^2-V(x)\), translation
symmetry in direction \(a\) gives \(J_a=m\dot x\cdot a\). If \(V\) is invariant
under all translations, linear momentum is conserved. Rotational symmetry about
axis \(a\) gives
\[
  J_a=(x\times m\dot x)\cdot a,
\]
so angular momentum is conserved when \(V\) is rotationally invariant.
\end{example}

\section{Basic Example: Kepler Motion From Central Forces}
\label{sec:kepler-lagrangian}

Kepler motion shows how cyclic coordinates, Noether's theorem, and reduction by
symmetry turn a force law into the geometry of conic sections.

\paragraph{Local notation.}
This section uses \(r\) first as a vector, \(r=x_1-x_2\), and then uses
\(\rho=\norm r\) for its length once the motion has been reduced to a plane.
The reduced mass is \(\mu=m_1m_2/(m_1+m_2)\), the coupling is
\(k=Gm_1m_2\), and \(\ell\) denotes the conserved angular momentum conjugate to
the polar angle \(\theta\).

The calculation proceeds in three reductions.  First the two-body problem is
written in the relative coordinate \(r=x_1-x_2\), so the center of mass is no
longer visible.  Second rotational invariance confines the motion to a fixed
plane perpendicular to the angular momentum.  Third the cyclic polar angle
\(\theta\) is replaced by its conserved momentum \(\ell\), leaving a
one-dimensional radial problem.  The conic-section formula is not guessed; it
is the result of these successive reductions.

\begin{assumption}[Kepler model]
We study a particle of reduced mass \(\mu\) moving in \(\R^3\setminus\{0\}\)
under an attractive inverse-square central potential. The Lagrangian is
\[
  L(r,\dot r)
  =
  \frac12\mu\norm{\dot r}^2+\frac{k}{\norm r},
  \qquad
  k=Gm_1m_2>0.
\]
The coordinate \(r=x_1-x_2\) is the relative separation of the two bodies.
\end{assumption}

Rotational invariance gives conservation of angular momentum
\[
  J=r\times p=\mu r\times\dot r.
\]
Since \(J\cdot r=0\) and \(J\cdot \dot r=0\), the motion stays in the plane
perpendicular to \(J\). In polar coordinates \((\rho,\theta)\) in that plane,
\[
  L=\frac12\mu(\dot\rho^2+\rho^2\dot\theta^2)+\frac{k}{\rho}.
\]
The angle \(\theta\) is cyclic, hence
\[
  \ell=\frac{\partial L}{\partial\dot\theta}=\mu\rho^2\dot\theta
\]
is conserved. The energy is
\[
  E=\frac12\mu\dot\rho^2+\frac{\ell^2}{2\mu\rho^2}-\frac{k}{\rho}.
\]
Thus the radial motion is one-dimensional motion in the effective potential
\[
  V_{\mathrm{eff}}(\rho)=\frac{\ell^2}{2\mu\rho^2}-\frac{k}{\rho}.
\]
Here \(\ell\) is the ordinary, massful angular momentum of the relative
two-body problem. Section~\ref{sec:detailed-central-forces} gives the
bookkeeping conversion to the specific angular momentum used in the asteroid
chapter.

\begin{figure}[htbp]
\centering
\begin{tikzpicture}[scale=1.0, >=Latex, font=\small]
  \draw[->] (0,0) -- (6.5,0) node[right] {\(\rho\)};
  \draw[->] (0,-1.9) -- (0,3.1) node[above] {\(V_{\mathrm{eff}}\)};
  \fill[blue!6] (0.624,-0.7) rectangle (2.52,-1.15);
  \draw[blue, thick, domain=0.36:6.1, samples=160]
    plot (\x,{1.1/(\x*\x)-2.2/\x});
  \draw[dashed] (0,-0.7) -- (6.2,-0.7);
  \node[fill=white, inner sep=1pt] (boundlabel) at (5.55,-1.34) {bound \(E<0\)};
  \draw[->, thin] (boundlabel.north west) -- (4.85,-0.7);
  \draw[dashed] (0,0.45) -- (6.2,0.45);
  \node[fill=white, inner sep=1pt] (scatlabel) at (4.85,0.92) {scattering \(E>0\)};
  \draw[->, thin] (scatlabel.south west) -- (4.2,0.45);
  \node[blue, fill=white, inner sep=1.5pt] at (4.75,-0.42)
    {\(V_{\mathrm{eff}}(\rho)\)};
  \draw[densely dotted] (1.0,-1.1) -- (1.0,-1.58);
  \fill (1.0,-1.1) circle (1.4pt);
  \node[fill=white, inner sep=1pt] (minlabel) at (1.9,-0.18) {minimum};
  \draw[->, thin] (minlabel.south west) -- (1.02,-1.05);
  \fill (0.624,-0.7) circle (1.3pt);
  \fill (2.52,-0.7) circle (1.3pt);
  \draw[<->, blue!70!black] (0.624,-1.58) -- (2.52,-1.58);
  \node[blue!70!black, fill=white, inner sep=1.5pt] at (1.70,-1.86)
    {allowed radial interval};
  \node[align=center, fill=white, inner sep=2pt] at (2.65,2.05)
    {centrifugal barrier\\near \(\rho=0\)};
\end{tikzpicture}
\caption{The Kepler problem reduces to one-dimensional radial motion in an
effective potential. For a fixed negative energy, the radial coordinate is
trapped between two turning points; for positive energy there is only a closest
approach before scattering.}
\label{fig:kepler-effective-potential}
\end{figure}

\begin{derivation}[Orbit equation]
Set \(u=1/\rho\). Since \(\ell=\mu\rho^2\dot\theta\),
\[
  \frac{\dd}{\dd t}=\dot\theta\frac{\dd}{\dd\theta}
  =
  \frac{\ell}{\mu}\,u^2\frac{\dd}{\dd\theta}.
\]
The radial Euler-Lagrange equation is
\[
  \mu(\ddot\rho-\rho\dot\theta^2)=-\frac{k}{\rho^2}.
\]
In terms of \(u(\theta)\) this becomes
\[
  \frac{\dd^2u}{\dd\theta^2}+u=\frac{\mu k}{\ell^2}.
\]
Therefore
\[
  u(\theta)=\frac{\mu k}{\ell^2}\left[1+e\cos(\theta-\theta_0)\right],
\]
equivalently
\[
  \boxed{
  \rho(\theta)=
  \frac{\ell^2/(\mu k)}{1+e\cos(\theta-\theta_0)}
  }.
\]
The parameter \(e\) is the eccentricity. The conic is an ellipse, parabola, or
hyperbola according as \(e<1\), \(e=1\), or \(e>1\).
\end{derivation}

The orbit equation uses \(\theta\) as the independent variable, so it describes
the shape of the orbit, not the rate at which the particle moves along it. The
time parametrization is recovered separately from
\(\ell=\mu\rho^2\dot\theta\).  This is a common pattern in mechanics:
symmetry first gives the geometry of a trajectory, while the conserved
quantity that generated the reduction then reconstructs the motion in time.

The Kepler problem also has the Laplace-Runge-Lenz vector
\[
  A=p\times J-\mu k\frac{r}{\norm r},
\]
which is conserved and points toward periapse. The conservation is derived in
Section~\ref{sec:detailed-central-forces}; here it is introduced as the extra
quantity that explains why the conic orientation does not precess in the ideal
inverse-square problem.

\begin{remark}[Bertrand's theorem and Kepler's hidden symmetry]
Bertrand's theorem (1873) gives the sharp structural statement: among central
force problems, only the inverse-square potential and the harmonic oscillator
make all bounded orbits closed. The Laplace-Runge-Lenz vector is the
inverse-square instance of this extra symmetry; in the negative-energy Kepler
problem it enlarges the rotational \(\SO(3)\) symmetry to the hidden
\(\SO(4)\) symmetry of bound motion. The vector itself goes back to Laplace's
celestial-mechanics work; Pauli (1926) used it in the hydrogen-atom
calculation, and Fock (1935) gave the momentum-sphere interpretation that makes
the \(\SO(4)\) explicit.
\end{remark}

\section{Basic Example: Rigid Body As A Lagrangian System}
\label{sec:rigid-body-lagrangian}

The rigid body is the first example in these notes whose configuration space is
a nonlinear Lie group. At the Lagrangian level, this already forces us to
distinguish coordinates, frames, angular velocity, and kinetic energy. Later
chapters return to the same system as a reduced Hamiltonian flow and as an
integrable model, but the essential geometric data are visible already here.
The reduced equations derived below were first written by Euler in 1758.

\paragraph{Local notation.}
Here \(R(t)\in\SO(3)\) maps body coordinates to space coordinates, \(X\) labels
a material point in the body frame, and \(x(t)=R(t)X\) is its position in
space.  The body angular velocity is \(\Omega\), the body angular momentum is
\(M=I\Omega\), and \(I\) is the inertia tensor.  The symbol \(I\) in this
section is therefore not an action variable.

\begin{assumption}[Free rigid body]
A material point \(X\) in the body frame is located at
\[
  x(t)=R(t)X,
  \qquad R(t)\in\SO(3),
\]
with the center of mass fixed at the origin. The body is rigid, so all dynamics
is in the attitude \(R(t)\). No external torque acts.
\end{assumption}

The body angular velocity is defined by
\[
  R^{-1}\dot R=\widehat\Omega,
  \qquad
  \widehat\Omega v=\Omega\times v.
\]
The kinetic energy is
\[
  L(R,\dot R)=T
  =
  \frac12\Omega^TI\Omega,
\]
where \(I\) is the inertia tensor in the body frame. In principal axes,
\[
  I=\diag(I_1,I_2,I_3),
  \qquad
  L=\frac12(I_1\Omega_1^2+I_2\Omega_2^2+I_3\Omega_3^2).
\]

\begin{derivation}[Euler equation from constrained variations]
A variation inside \(\SO(3)\) has the form
\[
  \delta R=R\widehat\Sigma,
\]
where \(\Sigma(t_0)=\Sigma(t_1)=0\). This is the left-trivialized form of a
tangent vector to \(\SO(3)\): multiplication by \(R^{-1}\) identifies
\(T_R\SO(3)\) with \(\mathfrak{so}(3)\). This implies
\[
  \delta\Omega=\dot\Sigma+\Omega\times\Sigma.
\]
Indeed, differentiating \(R^TR=I\) gives
\[
  R^T\delta R+(\delta R)^TR=0,
\]
so \(R^T\delta R\in\mathfrak{so}(3)\), which is why
\(\delta R=R\widehat\Sigma\). Then
\[
  \delta(R^T\dot R)
  =
  \delta R^T\dot R+R^T\delta\dot R
  =
  -\widehat\Sigma\widehat\Omega
  +\widehat{\dot\Sigma}+\widehat\Omega\widehat\Sigma
  =
  \widehat{\dot\Sigma}+[\widehat\Omega,\widehat\Sigma],
\]
and the hat map satisfies
\[
  [\widehat\Omega,\widehat\Sigma]=\widehat{\Omega\times\Sigma}.
\]
On a test vector \(v\), this identity is
\[
  [\widehat\Omega,\widehat\Sigma]v
  =
  \Omega\times(\Sigma\times v)-\Sigma\times(\Omega\times v)
  =
  (\Omega\times\Sigma)\times v,
\]
which is the Jacobi identity for the cross product written as a commutator.
This is the stated formula under the hat isomorphism.
Let \(M=\partial L/\partial\Omega=I\Omega\). Then
\[
\begin{aligned}
  \delta S
  &=
  \int M\cdot(\dot\Sigma+\Omega\times\Sigma)\,\dd t\\
  &=
  \int(-\dot M+M\times\Omega)\cdot\Sigma\,\dd t.
\end{aligned}
\]
Stationarity gives
\[
  \dot M=M\times\Omega,
\]
or, with the hat convention \(\widehat\Omega v=\Omega\times v\) used
throughout,
\[
  \dot M+\Omega\times M=0.
\]
In principal axes this is
\[
\begin{aligned}
  I_1\dot\Omega_1 &= (I_2-I_3)\Omega_2\Omega_3,\\
  I_2\dot\Omega_2 &= (I_3-I_1)\Omega_3\Omega_1,\\
  I_3\dot\Omega_3 &= (I_1-I_2)\Omega_1\Omega_2.
\end{aligned}
\]
\end{derivation}

\section{What To Take Forward}
\label{sec:lagrangian-take-forward}

\begin{itemize}
\item Hamilton's principle starts from a configuration manifold \(Q\), paths
\(q(t)\in Q\), an action \(S=\int L\,\dd t\), and a chosen class of admissible
variations. For the fixed-endpoint boundary-value problem, the variational
postulate requires stationarity under endpoint-fixing variations. That
stationarity gives the Euler-Lagrange equations, and the endpoint term defines
conjugate momentum when endpoint variations are allowed.
\item Natural Lagrangians encode kinetic energy as a mass metric. Forces,
constraints, and nonconservative virtual work enter as covectors, multipliers,
or potentials, depending on how they act on virtual displacements.
\item Noether's theorem turns continuous symmetries into conserved quantities.
Kepler motion and the rigid body show how these quantities organize exact
solutions and prepare the Legendre transform to Hamiltonian phase space.
\end{itemize}

\section{Chapter-End Laboratories And Exercises}
\label{sec:lagrangian-chapter-end-labs}

The laboratories below test the Lagrangian ideas of this chapter: coordinate
independence, symmetry reduction, central-force stability, time-dependent
Lagrangians, effective potentials, and contact forces.

\subsection{Form-Invariance Of Lagrange Equations}
\label{sec:lab-point-transformations}

We begin with the coordinate-independence of the Euler-Lagrange equations. Let
\[
  q^i=q^i(s^1,\ldots,s^n,t)
\]
be a regular time-dependent point transformation. Regular means that the
Jacobian \(\partial q^i/\partial s^a\) has rank \(n\). The transformed
Lagrangian is
\[
  L'(s,\dot s,t)=L(q(s,t),\dot q(s,\dot s,t),t),
\]
where
\[
  \dot q^i
  =
  \frac{\partial q^i}{\partial s^a}\dot s^a
  +
  \frac{\partial q^i}{\partial t}.
\]

\paragraph{Local notation.}
The symbols \(q^i\) and \(s^a\) are two coordinate systems on the same
configuration manifold, not two different physical systems.  The index \(a\)
labels the new coordinates, and regularity means the coordinate change has an
invertible Jacobian at each point under discussion.

Two elementary identities do the work:
\[
  \frac{\partial \dot q^i}{\partial \dot s^a}
  =
  \frac{\partial q^i}{\partial s^a},
  \qquad
  \frac{\dd}{\dd t}
  \left(
    \frac{\partial q^i}{\partial s^a}
  \right)
  =
  \frac{\partial \dot q^i}{\partial s^a}.
\]
Using the chain rule,
\[
  \frac{\partial L'}{\partial \dot s^a}
  =
  \frac{\partial L}{\partial \dot q^i}
  \frac{\partial q^i}{\partial s^a},
\]
and
\[
  \frac{\partial L'}{\partial s^a}
  =
  \frac{\partial L}{\partial q^i}
  \frac{\partial q^i}{\partial s^a}
  +
  \frac{\partial L}{\partial \dot q^i}
  \frac{\partial \dot q^i}{\partial s^a}.
\]
Therefore
\[
\begin{aligned}
  \frac{\dd}{\dd t}\frac{\partial L'}{\partial \dot s^a}
  -
  \frac{\partial L'}{\partial s^a}
  &=
  \left(
    \frac{\dd}{\dd t}\frac{\partial L}{\partial\dot q^i}
    -
    \frac{\partial L}{\partial q^i}
  \right)
  \frac{\partial q^i}{\partial s^a}.
\end{aligned}
\]
If \(q(t)\) satisfies the Euler-Lagrange equations, the right side vanishes;
if the transformation is regular, the converse holds as well. Thus the
Euler-Lagrange equation is not tied to a preferred coordinate system.

\subsection{Central Forces From First Principles}
\label{sec:detailed-central-forces}

The central-force problem is the prototype of reduction by symmetry. Rotational
invariance gives conservation of angular momentum; angular momentum confines
the motion to a plane; polar coordinates in that plane reduce the dynamics to a
one-dimensional radial equation with an effective potential. The later Kepler
and precession calculations are refinements of this same reduction.

The lesson is broader than celestial mechanics. A symmetry gives a conserved
quantity; the conserved quantity can remove a variable; the reduced equation
then describes motion on a smaller space. Central forces are a transparent first
example because the geometry is visible in ordinary Euclidean space.

Let a particle of mass \(m\) move in a central potential \(V(r)\), where
\[
  r=\norm{x},
  \qquad
  x\in\R^3\setminus\{0\}.
\]

\paragraph{Local notation.}
In this laboratory \(r=\norm{x}\) is the radial distance, \(J=x\times p\) is
the angular-momentum vector, and \(\ell=\norm J\) is its scalar value in the
orbital plane.  The prime in \(V'(r)\) denotes differentiation with respect to
the radial coordinate \(r\); later, a prime on \(u(\theta)\) denotes
\(\dd/\dd\theta\).

The Lagrangian is
\[
  L=\frac12m\norm{\dot x}^2-V(r).
\]
Rotational invariance gives angular momentum
\[
  J=x\times p=m x\times\dot x.
\]
Since
\[
  \frac{\dd J}{\dd t}
  =
  \dot x\times m\dot x+x\times m\ddot x
  =
  x\times\left(-V'(r)\frac{x}{r}\right)=0,
\]
the vector \(J\) is constant. Because \(J\cdot x=0\) and
\(J\cdot\dot x=0\), the motion lies in the plane perpendicular to \(J\). Use
polar coordinates \((r,\theta)\) in this plane:
\[
  L=\frac12m(\dot r^2+r^2\dot\theta^2)-V(r).
\]
The angle \(\theta\) is cyclic, so
\[
  \ell=\frac{\partial L}{\partial\dot\theta}=mr^2\dot\theta
\]
is conserved. The energy is
\[
  E=\frac12m\dot r^2+\frac{\ell^2}{2mr^2}+V(r).
\]
Thus the radial motion is one-dimensional motion in
\[
  V_{\mathrm{eff}}(r)
  =
  V(r)+\frac{\ell^2}{2mr^2}.
\]
Here \(\ell\) is the massful angular momentum. To compare with the
specific-energy Kepler convention used in
Section~\ref{sec:kepler-perturbation}, divide the Hamiltonian, energy, and
angular momentum by the particle mass; then
\(\ell_{\mathrm{spec}}=\ell/m=r^2\dot\theta\). The orbit equation is unchanged
after this bookkeeping conversion, provided \(k/m\) is read as the
gravitational parameter \(GM\).

\subsubsection{Binet Equation}
\label{subsec:detailed-binet}

Let
\[
  u(\theta)=\frac1r.
\]

\paragraph{Local notation.}
From this point until the end of the Binet derivation, a prime means
\(\dd/\dd\theta\), not a derivative with respect to \(r\) or time.  The function
\(u\) has units of inverse length and is introduced only to write the orbit
shape as a function of the polar angle.

Since
\[
  \dot\theta=\frac{\ell}{mr^2}=\frac{\ell}{m}u^2,
\]
we have
\[
  \frac{\dd}{\dd t}
  =
  \frac{\ell}{m}u^2\frac{\dd}{\dd\theta}.
\]
Also
\[
  \dot r
  =
  \frac{\dd}{\dd t}(u^{-1})
  =
  -u^{-2}\dot u
  =
  -\frac{\ell}{m}u',
\]
where a prime denotes \(\dd/\dd\theta\). Differentiating once more,
\[
  \ddot r
  =
  -\frac{\ell}{m}\frac{\dd u'}{\dd t}
  =
  -\frac{\ell^2}{m^2}u^2u''.
\]
The radial equation is
\[
  m(\ddot r-r\dot\theta^2)=F(r),
  \qquad
  F(r)=-V'(r).
\]
Substitute the formulas above:
\[
  m\left(
    -\frac{\ell^2}{m^2}u^2u''
    -
    \frac1u\frac{\ell^2}{m^2}u^4
  \right)
  =
  -\frac{\ell^2}{m}u^2(u''+u)
  =
  F(1/u).
\]
Therefore
\[
  \boxed{
  u''+u
  =
  -\frac{m}{\ell^2u^2}F(1/u).
  }
\]
This is Binet's equation. It turns a central-force problem into a differential
equation for the orbit shape \(r(\theta)\).

\subsubsection{Kepler Conics And The Laplace-Runge-Lenz Vector}
\label{subsec:detailed-kepler-lrl}

For the Kepler potential
\[
  V(r)=-\frac{k}{r},
  \qquad
  k>0,
\]
the force is
\[
  F(r)=-\frac{k}{r^2}.
\]
Binet's equation becomes
\[
  u''+u=\frac{mk}{\ell^2}.
\]
The solution is
\[
  u(\theta)=\frac{mk}{\ell^2}\left[1+e\cos(\theta-\theta_0)\right],
\]
equivalently
\[
  r(\theta)=\frac{\ell^2/(mk)}
  {1+e\cos(\theta-\theta_0)}.
\]
This is a conic section with focus at the force center.

The extra conserved vector is
\[
  A=p\times J-mk\frac{x}{r}.
\]
To verify conservation, write \(\hat r=x/r\). Since \(J\) is constant,
\[
  \dot A=\dot p\times J-mk\dot{\hat r}.
\]
The force law gives
\[
  \dot p=-\frac{k}{r^2}\hat r.
\]
Use
\[
  J=mx\times\dot x=mr^2\dot\theta\,\hat z=\ell\hat z
\]
in the orbital plane. Then
\[
  \dot{\hat r}=\dot\theta\,\hat\theta,
\]
and
\[
  \dot p\times J
  =
  -\frac{k}{r^2}\hat r\times(\ell\hat z)
  =
  \frac{k\ell}{r^2}\hat\theta
  =
  mk\dot\theta\,\hat\theta
  =
  mk\dot{\hat r}.
\]
Hence
\[
  \dot A=0.
\]
The vector \(A\) points toward periapse. Taking the dot product with \(x\),
\[
  A\cdot x
  =
  (p\times J)\cdot x-mkr.
\]
By cyclic permutation,
\[
  (p\times J)\cdot x=J\cdot(x\times p)=J\cdot J=\ell^2.
\]
Therefore
\[
  A r\cos(\theta-\theta_0)=\ell^2-mkr,
\]
which rearranges to
\[
  r=\frac{\ell^2/(mk)}{1+(A/(mk))\cos(\theta-\theta_0)}.
\]
Thus the eccentricity is
\[
  e=\frac{\norm A}{mk}.
\]
The conic equation and the hidden conservation law are the same fact written
in two languages.

\subsubsection[Small Periapse Precession]{Small Periapse Precession From A
\texorpdfstring{\(1/r^3\)}{1/r^3} Perturbation}
\label{subsec:detailed-precession}

A useful perturbation of the Kepler potential is
\[
  V(r)=-\frac{k}{r}-\frac{\epsilon}{r^3},
  \qquad
  0<|\epsilon|\ll1.
\]
This is a pedagogical proxy for any small correction that breaks the exact
inverse-square symmetry while remaining central. The Schwarzschild laboratory
below gives the relativistic source of a closely related leading precession
effect.
The radial force is
\[
  F(r)=-V'(r)=-\frac{k}{r^2}-\frac{3\epsilon}{r^4}.
\]
Binet's equation gives
\[
  u''+u
  =
  \frac{mk}{\ell^2}
  +
  \frac{3m\epsilon}{\ell^2}u^2.
\]
For a nearly circular orbit, write
\[
  u(\theta)=u_0+\eta(\theta),
  \qquad
  |\eta|\ll u_0.
\]
The constant \(u_0\) solves
\[
  u_0=\frac{mk}{\ell^2}+\frac{3m\epsilon}{\ell^2}u_0^2.
\]
Linearizing in \(\eta\),
\[
  \eta''+\eta
  =
  \frac{6m\epsilon u_0}{\ell^2}\eta,
\]
where terms of order \(\epsilon\eta^2\) and higher are discarded. The small
quantity \(\eta/u_0\) measures the radial oscillation amplitude, while
\(\epsilon\) measures the departure from the Kepler force law; the displayed
linear equation keeps the terms needed for the leading frequency shift.
Equivalently,
\[
  \eta''+\kappa^2\eta=0,
  \qquad
  \kappa^2=1-\frac{6m\epsilon u_0}{\ell^2}.
\]
The radial oscillation has angular period
\[
  \Delta\theta_r=\frac{2\pi}{\kappa}.
\]
The periapse advance per radial cycle is the excess over \(2\pi\):
\[
  \Delta\varpi
  =
  2\pi\left(\frac1\kappa-1\right)
  \approx
  \frac{6\pi m\epsilon u_0}{\ell^2}.
\]
Thus the \(1/r^3\) correction breaks the closed Kepler ellipse and produces
precession. The same algebraic structure appears in general relativity.
Einstein used it in 1915 to explain Mercury's anomalous perihelion advance; the
present subsection treats only a classical effective potential, but the
calculation has the same shape.

\subsection{Damped Oscillator From A Time-Dependent Lagrangian}
\label{sec:lab-damped-oscillator}

Let
\[
  L(q,\dot q,t)
  =
  e^{\gamma t}
  \left(
    \frac12m\dot q^2-\frac12kq^2
  \right).
\]

\paragraph{Local notation.}
The coordinate \(q(t)\) is the physical oscillator displacement, \(s(t)\) is a
rescaled variable introduced for the calculation, \(\gamma\) is the damping
rate, \(m\) is the mass, and \(k\) is the spring constant.  The explicit factor
\(e^{\gamma t}\) makes the Lagrangian time-dependent.

The Euler-Lagrange equation is
\[
  \frac{\dd}{\dd t}(e^{\gamma t}m\dot q)+e^{\gamma t}kq=0,
\]
equivalently
\[
  \boxed{
  m\ddot q+m\gamma\dot q+kq=0.
  }
\]
Although this is a dissipative equation, it comes from a time-dependent
Lagrangian. The canonical energy is not conserved because \(L\) depends
explicitly on time.

Set
\[
  s=e^{\gamma t/2}q,
  \qquad
  q=e^{-\gamma t/2}s.
\]
Then
\[
  \dot q=e^{-\gamma t/2}\left(\dot s-\frac{\gamma}{2}s\right).
\]
Substituting into the action gives
\[
  L
  =
  \frac12m\left(\dot s-\frac{\gamma}{2}s\right)^2
  -
  \frac12ks^2.
\]
The cross term is a total derivative:
\[
  -\frac12m\gamma s\dot s
  =
  -\frac{\dd}{\dd t}\left(\frac14m\gamma s^2\right).
\]
Discarding that total derivative, the equivalent Lagrangian is
\[
  L_{\mathrm{eff}}
  =
  \frac12m\dot s^2
  -
  \frac12m
  \left(
    \frac{k}{m}-\frac{\gamma^2}{4}
  \right)s^2.
\]
Thus
\[
  \ddot s+\left(\frac{k}{m}-\frac{\gamma^2}{4}\right)s=0.
\]
The transformation has moved damping into a shifted oscillator frequency.

\begin{figure}[htbp]
\centering
\begin{tikzpicture}[x=0.82cm,y=1.0cm, >=Latex, font=\small]
  \draw[->] (-0.2,0) -- (8.45,0) node[right] {\(t\)};
  \draw[->] (0,-1.45) -- (0,1.6) node[above] {amplitude};
  \draw[blue!70!black, dashed, domain=0:8.0, samples=150]
    plot (\x,{1.10*exp(-0.25*\x)});
  \draw[blue!70!black, dashed, domain=0:8.0, samples=150]
    plot (\x,{-1.10*exp(-0.25*\x)});
  \draw[blue!70!black, very thick, domain=0:8.0, samples=260]
    plot (\x,{1.10*exp(-0.25*\x)*cos(3.25*\x r)});
  \draw[red!75!black, dashed] (0,1.10) -- (8.0,1.10);
  \draw[red!75!black, dashed] (0,-1.10) -- (8.0,-1.10);
  \draw[red!75!black, thick, domain=0:8.0, samples=260]
    plot (\x,{1.10*cos(3.25*\x r)});
  \node[blue!70!black, fill=white, inner sep=1.5pt] (qlabel) at (5.35,-0.92) {\(q(t)\)};
  \draw[blue!70!black, ->] (qlabel.north) -- (5.25,-0.18);
  \node[red!75!black, fill=white, inner sep=1.5pt] at (5.85,1.24)
    {\(s=e^{\gamma t/2}q\)};
  \node[red!75!black, fill=white, inner sep=1pt] at (1.18,1.24)
    {constant envelope};
  \node[align=center, fill=white, inner sep=1.5pt] at (4.1,-1.78)
    {rescaling absorbs the exponential damping envelope};
\end{tikzpicture}
\caption{Damped-oscillator laboratory. In the original coordinate \(q\), the
motion has a decaying envelope. The rescaled coordinate \(s=e^{\gamma t/2}q\)
obeys an undamped oscillator equation with shifted frequency.}
\label{fig:lab-damped-oscillator-rescaling}
\end{figure}
\FloatBarrier

\subsection{Rotating Hoop: Effective Potential And Bifurcation}
\label{sec:lab-rotating-hoop}

A bead of mass \(m\) moves without friction on a hoop of radius \(a\). The hoop
is vertical and rotates about its vertical diameter with imposed angular speed
\(\omega\). Let \(\theta\) be the angle of the bead measured from the downward
vertical. The bead position has speed
\[
  v^2=a^2\dot\theta^2+a^2\omega^2\sin^2\theta.
\]

\paragraph{Local notation.}
The angular speed \(\omega\) is imposed by an external motor and is not the
bead's generalized velocity.  The only dynamical coordinate is \(\theta(t)\);
the term \(a\omega\sin\theta\) is the speed of the bead due to the hoop's
rotation about the vertical axis.

The Lagrangian is
\[
  L
  =
  \frac12ma^2\dot\theta^2
  +
  \frac12ma^2\omega^2\sin^2\theta
  +
  mga\cos\theta.
\]

\begin{figure}[htbp]
\centering
\begin{tikzpicture}[scale=1.0, >=Latex, font=\small]
  \begin{scope}
    \draw[->] (0,-2.05) -- (0,2.25) node[above] {rotation axis};
    \draw[dashed] (0,-1.85) -- (0,1.85);
    \draw[thick] (0,0) circle (1.65);
    \draw[fill=red!70!black] (1.17,-1.17) circle (0.09);
    \draw[thick] (0,0) -- (1.17,-1.17);
    \draw[->] (0,-0.75) arc (-90:-45:0.75);
    \node[fill=white, inner sep=1pt] at (0.55,-0.95) {\(\theta\)};
    \draw[->, blue!70!black] (1.17,-1.17) -- (1.66,-1.17)
      node[pos=0.75, below=2pt, fill=white, inner sep=1pt]
      {\(a\omega\sin\theta\)};
    \node[below, fill=white, inner sep=1pt] at (0,-1.72) {downward vertical};
    \node[fill=white, inner sep=1pt] at (0,-2.40) {bead on rotating vertical hoop};
  \end{scope}
  \begin{scope}[xshift=5.0cm, xscale=0.58]
    \draw[->] (-3.45,0) -- (3.45,0) node[right] {\(\theta\)};
    \draw[->] (0,-1.5) -- (0,1.55) node[above] {\(V_{\mathrm{eff}}\)};
    \draw[gray!60] (-3.1416,-0.05) -- (-3.1416,0.05)
      node[below=3pt] {\(-\pi\)};
    \draw[gray!60] (3.1416,-0.05) -- (3.1416,0.05)
      node[below=3pt] {\(\pi\)};
    \draw[blue!70!black, thick, domain=-3.1416:3.1416, samples=180]
      plot (\x,{-0.35*sin(\x r)*sin(\x r)-0.80*cos(\x r)});
    \draw[red!75!black, thick, domain=-3.1416:3.1416, samples=180]
      plot (\x,{-1.45*sin(\x r)*sin(\x r)-0.25*cos(\x r)+0.48});
    \draw[blue!70!black, thick] (0.62,1.22) -- (1.02,1.22);
    \node[blue!70!black, align=left, fill=white, inner sep=1.5pt] at (2.02,1.22)
      {\(\omega<\omega_c\): one minimum};
    \draw[red!75!black, thick] (0.62,0.84) -- (1.02,0.84);
    \node[red!75!black, align=left, fill=white, inner sep=1.5pt] at (2.02,0.84)
      {\(\omega>\omega_c\): two minima};
  \end{scope}
\end{tikzpicture}
\caption{Rotating-hoop laboratory. The geometry fixes the kinetic term
\(a^2\omega^2\sin^2\theta\), and the effective potential changes from one
bottom minimum to two off-bottom minima when \(\omega\) crosses
\(\omega_c=\sqrt{g/a}\). The potential curves are schematic and emphasize the
bifurcation, not numerical scale.}
\label{fig:lab-rotating-hoop}
\end{figure}
\FloatBarrier

This is a one-dimensional system with effective potential
\[
  V_{\mathrm{eff}}(\theta)
  =
  -\frac12ma^2\omega^2\sin^2\theta
  -
  mga\cos\theta.
\]
Equilibria satisfy
\[
  \frac{\dd V_{\mathrm{eff}}}{\dd\theta}
  =
  ma\sin\theta(g-a\omega^2\cos\theta)=0.
\]
Thus either
\[
  \theta=0,\pi,
\]
or else
\[
  \cos\theta=\frac{g}{a\omega^2}.
\]
The off-bottom equilibria exist only when
\[
  \omega\geq\omega_c=\sqrt{\frac{g}{a}}.
\]
At \(\omega=\omega_c\), the bottom equilibrium changes stability and two
symmetric equilibria appear. This is a concrete mechanics realization of a
pitchfork bifurcation in an effective potential.

\subsection{Rolling Contact: Contact Forces As Multipliers}
\label{sec:lab-rolling-contact}

Rolling examples force one to distinguish coordinates, constraints, and
reaction forces.

\paragraph{Local notation.}
The symbol \(N\) below denotes the normal contact force, not the dimension of a
configuration space.  The angles \(\theta\) and \(\phi\) describe orbital
motion of the center and spin of the rolling body, respectively; the rolling
condition relates their time derivatives.

\subsubsection{Hoop Rolling On A Fixed Cylinder}

A hoop of mass \(m\) and radius \(r\) rolls without slipping on the outside of a
fixed cylinder of radius \(R\). Let \(\theta\) be the angular position of the
hoop center measured from the top. The center of mass moves on a circle of
radius \(R+r\). The rolling condition gives
\[
  r\dot\phi=(R+r)\dot\theta,
\]
where \(\phi\) is the spin angle of the hoop. Since the hoop has
\[
  I_{\mathrm{hoop}}=mr^2,
\]
the kinetic energy is
\[
  T=
  \frac12m(R+r)^2\dot\theta^2
  +
  \frac12mr^2\dot\phi^2
  =
  m(R+r)^2\dot\theta^2.
\]
Starting from rest at the top, energy gives
\[
  m(R+r)^2\dot\theta^2
  =
  mg(R+r)(1-\cos\theta).
\]
The radial equation is
\[
  mg\cos\theta-N=m(R+r)\dot\theta^2,
\]
where \(N\) is the normal force. Loss of contact occurs at \(N=0\). Using the
energy relation at \(N=0\),
\[
  g\cos\theta=(R+r)\dot\theta^2
  =
  g(1-\cos\theta),
\]
so
\[
  \boxed{\cos\theta=\frac12.}
\]
The hoop leaves the cylinder at \(\theta=\pi/3\) from the top.

\begin{figure}[htbp]
\centering
\begin{tikzpicture}[scale=1.0, >=Latex, font=\small]
  \draw[thick, fill=gray!8] (0,0) circle (1.4);
  \coordinate (cyl) at (0,0);
  \coordinate (hoop) at (1.65,1.65);
  \coordinate (contact) at (0.99,0.99);
  \draw[thick, fill=blue!7] (hoop) circle (0.45);
  \draw[dashed] (cyl) -- (hoop);
  \draw[dashed] (0,0) -- (0,2.35);
  \draw[->] (0,1.9) arc (90:45:1.9);
  \node[fill=white, inner sep=1pt] at (0.78,2.15) {\(\theta\)};
  \fill[black] (contact) circle (1.1pt);
  \draw[black!60, densely dashed] (0.70,1.28) -- (1.28,0.70);
  \node[fill=white, inner sep=1pt] at (0.58,0.70) {tangent};
  \draw[->, red!75!black, very thick] (hoop) -- (2.35,2.35);
  \node[red!75!black, fill=white, inner sep=1pt] at (2.55,2.45) {\(N\)};
  \draw[->, black!70, thick] (hoop) -- (1.65,0.75)
    node[below] {\(mg\)};
  \draw[->, blue!70!black, thick] (2.1,1.65) arc (0:65:0.45);
  \node[blue!70!black, fill=white, inner sep=1pt] at (2.36,1.42) {\(\dot\phi\)};
  \node[align=center] at (0,-1.95)
    {loss of contact occurs when the\\normal force reaches \(N=0\)};
  \node[fill=white, inner sep=1.5pt] at (-1.08,-1.12) {fixed cylinder \(R\)};
  \node[fill=white, inner sep=1.5pt] at (2.78,2.58) {rolling hoop \(r\)};
\end{tikzpicture}
\caption{Rolling-contact laboratory. The radial force balance must be kept
until the end because the contact force \(N\) is the diagnostic for leaving the
surface. Eliminating the rolling constraint gives the motion, but \(N=0\)
decides the loss-of-contact angle.}
\label{fig:lab-rolling-contact}
\end{figure}
\FloatBarrier

\subsubsection{Solid Cylinder Inside A Hollow Cylinder}

Now consider a solid cylinder of radius \(R\) rolling inside a fixed hollow
cylinder of radius \(2R\). The center of mass moves on a circle of radius
\(R\). With rolling imposed, the kinetic energy is
\[
  T
  =
  \frac12MR^2\dot\theta^2
  +
  \frac14MR^2\dot\theta^2,
\]
because a solid cylinder has moment of inertia \(I=(1/2)MR^2\). Taking the
potential as \(V=-MgR\cos\theta\), the Lagrangian is
\[
  L=\frac34MR^2\dot\theta^2+MgR\cos\theta.
\]
The equation of motion is
\[
  \frac32MR^2\ddot\theta+MgR\sin\theta=0,
\]
equivalently
\[
  \frac32\ddot\theta+\frac{g}{R}\sin\theta=0.
\]
For small oscillations,
\[
  \boxed{
  \omega_{\mathrm{small}}=\sqrt{\frac{2g}{3R}}.
  }
\]

\subsection{Schwarzschild-Type Central-Force Stability}
\label{sec:lab-schwarzschild-stability}

Consider the central Lagrangian
\[
  L=-m\sqrt{
    A(r)-\frac{\dot r^2}{A(r)}-r^2\dot\theta^2
  },
  \qquad
  A(r)=1-\frac{R}{r},
  \qquad r>R.
\]

\paragraph{Local notation.}
In this laboratory \(R\) is the Schwarzschild radius parameter, while \(r\) is
the radial coordinate of the orbit.  The quantity \(L_z\) is angular momentum
about the symmetry axis; it should not be confused with the Lagrangian \(L\).
The symbol \(B\) is only an abbreviation for the expression inside the square
root.

This is the geodesic Lagrangian for the Schwarzschild metric, derived by
Schwarzschild in 1916, with horizon at \(r=R\), restricted to the equatorial
plane, in units \(c=1\). Here the parameter is Schwarzschild coordinate time.
The model is therefore the reduced square-root Lagrangian for spatial motion,
not the reparametrization-invariant proper-time action. We use it as a
classical central-force laboratory whose effective-potential analysis already
yields the correct innermost stable circular-orbit radius.
Define
\[
  B=A-\frac{\dot r^2}{A}-r^2\dot\theta^2.
\]
The conserved energy and angular momentum are
\[
  E=\frac{mA}{\sqrt B},
  \qquad
  L_z=\frac{mr^2\dot\theta}{\sqrt B}.
\]
The angular momentum follows from the cyclic coordinate \(\theta\). The energy
is the autonomous-Lagrangian first integral traditionally called the Beltrami
integral. It is conserved here because this reduced Lagrangian has no explicit
time dependence:
\[
  \dot r\,\frac{\partial L}{\partial\dot r}
  +
  \dot\theta\,\frac{\partial L}{\partial\dot\theta}
  -
  L
  =
  \frac{mA}{\sqrt B}.
\]
Equivalently, it is the time-translation Noether charge
\(\dot q^i\partial L/\partial\dot q^i-L\); after the Legendre transform the
same expression is the Hamiltonian.
Eliminating \(\dot\theta\) and \(B\) gives the radial equation in effective
potential form:
\[
  \dot r^2
  =
  \frac{A^2}{E^2}
  \left[
    E^2-A\left(m^2+\frac{L_z^2}{r^2}\right)
  \right].
\]
Thus circular orbits are critical points of
\[
  U_{\mathrm{eff}}(r)
  =
  A(r)\left(m^2+\frac{L_z^2}{r^2}\right)
\]
at fixed \(L_z\). The circular-orbit condition \(U_{\mathrm{eff}}'(r_0)=0\)
gives
\[
  L_z^2
  =
  \frac{Rm^2r_0^2}{2r_0-3R}.
\]
Hence circular orbits require \(r_0>3R/2\). Stability is determined by
\[
  U_{\mathrm{eff}}''(r_0)
  =
  \frac{2Rm^2(r_0-3R)}
  {r_0^3(2r_0-3R)}.
\]
Therefore
\[
  \boxed{
  r_0>3R \text{ is stable},\qquad
  \frac32R<r_0<3R \text{ is unstable}.
  }
\]
This is the same structural calculation used in more realistic relativistic
orbit problems: circular-orbit stability is a second-derivative question for
an effective potential at fixed angular momentum.

\begin{figure}[htbp]
\centering
\begin{tikzpicture}[x=0.62cm,y=1.25cm, >=Latex, font=\small]
  \draw[->] (0,0) -- (10.55,0) node[right] {\(r/R\)};
  \draw[->] (0,-0.25) -- (0,2.25) node[above] {\(U_{\mathrm{eff}}\)};
  \draw[dashed] (1.0,-0.15) -- (1.0,2.05);
  \draw[densely dashed, red!55!black] (2.0,-0.15) -- (2.0,1.62);
  \draw[dashed] (3.0,-0.15) -- (3.0,2.05);
  \draw[densely dashed, blue!55!black] (6.0,-0.15) -- (6.0,1.10);
  \node[below, fill=white, inner sep=1pt] at (1.0,-0.15) {\(R\)};
  \node[below, fill=white, inner sep=1pt] at (2.0,-0.15) {\(2R\)};
  \node[below, fill=white, inner sep=1pt] at (3.0,-0.15) {\(3R\)};
  \node[below, fill=white, inner sep=1pt] at (6.0,-0.15) {\(6R\)};
  \draw[red!75!black, thick, domain=1.25:10.0, samples=260]
    plot (\x,{1.50+5.5*(-1/\x+4/(\x*\x)-4/(\x*\x*\x))});
  \fill[red!75!black] (2.0,1.50) circle (1.3pt);
  \node[red!75!black, fill=white, inner sep=1.5pt] at (3.75,2.02)
    {unstable circular orbit};
  \fill[blue!70!black] (6.0,1.09) circle (1.3pt);
  \node[blue!70!black, fill=white, inner sep=1.5pt] at (7.55,0.82)
    {stable circular orbit};
  \draw[blue!70!black, thick, ->] (7.15,0.88) -- (6.12,1.06);
  \draw[red!75!black, thick, ->] (3.25,1.88) -- (2.05,1.52);
  \node[align=center, fill=white, inner sep=1.5pt] at (5.3,-0.72)
    {stability is the sign of \(U_{\mathrm{eff}}''(r_0)\)\\
     at fixed angular momentum};
\end{tikzpicture}
\caption{Effective-potential picture for the Schwarzschild-type laboratory.
Circular orbits are critical points. The plotted curve is an affine rescaling of
\(-1/r+4/r^2-4/r^3\), whose circular-orbit extrema occur at \(r=2R\) and
\(r=6R\): the inner extremum is unstable and the outer one is stable.}
\label{fig:lab-schwarzschild-effective-potential}
\end{figure}
\FloatBarrier

\section{Associated Demos}
\label{sec:lagrangian-associated-demos}

This chapter is mostly analytic, but two later demos are natural companions for
its examples:
\begin{itemize}
\item \path{demos/python/hamiltonian_pendulum.py}, read after the Legendre
transform in Chapter~\ref{chap:hamiltonian-phase-space}, turns effective
potential thinking into a phase portrait.
\item \path{demos/python/rigid_body_euler_top.py}, read after
Chapter~\ref{chap:rigid-integrability}, continues the Lagrangian rigid-body
derivation into reduced Euler equations and attitude reconstruction.
\end{itemize}
The same examples are first derived from the action, then revisited in
phase-space or reduced variables. The computation adds a later layer to the same
model.
